% NUMBERS. Simulation results come from sc-workshop-paper/results_tables/*.md,
% regenerated by scripts/make_results_tables.py; the prose cites cells, not
% hand-copied constants, so swapping simulation for measurement means re-running
% that script and re-reading the same cells:
%
%   Stall Reduction     -> 02_budget_sweep.md
%   End-to-End Runtime  -> 02_budget_sweep.md, 03_compute_sweep.md
%   DMSH comparison     -> results/eval_q1_q4/eval_q1_summary.csv (MEASURED),
%                          faceted to gpu_name == "NVIDIA L40S"; stall classes
%                          from eval_stall_taxonomy.csv. The ONLY measured
%                          end-to-end subsection. Do not let it be pooled
%                          across accelerator generations.
%   Ablations           -> 01_attribution_ladder.md, 06_prefetch_variants.md
%   CPU interference    -> results/bench_preactivation_interference.json and
%                          results/bench_potential_activation_rusage_*.json
%                          (both MEASURED)
%   Sensitivity         -> 07_objective_check.md, 02_budget_sweep.md,
%                          04_accuracy_sweep.md
%
% OPEN ITEM: the exact-selection column of 04_accuracy_sweep.md runs deeply
% negative at 560 GB / 2 slots. The system specifies a predictor reporting "not
% within H" rather than "never"; those cells appear to have been generated
% WITHOUT that bounded-lookahead reporting, so they may measure an arm the
% design does not describe. The exact arm is omitted from the prose below until
% that is resolved.

\section{Evaluation}

\subsection{Setup}

Results are trace-driven simulation over measured resource constants. Park
footprints, cold-boot and wake costs, activated sizes and per-format
construction rates are measured on hardware; schedules and the clock are
synthetic. Each reported value averages 3 popularity orderings $\times$ 2
schedule populations $\times$ 10 schedules of 48 resource needs. The
comparison against a byte-oriented tiering system
(Section~\ref{sec:dmsh}) is measured end-to-end and is marked as such.

We report against two baselines. \textsc{Never} retains nothing: a displaced
model is terminated and a data artifact is reconstructed on every use. It is
the honest denominator for the mechanism. \textsc{Lru} retains both classes
but ranks by recency alone; it is already a system, so it is the line any
policy claim must clear. Reporting only against \textsc{Never} would credit
arbitration for benefits that plain retention already delivers.

\subsection{Stall Reduction}

At the production 256~GB allocation with a single device slot, \textsc{Lru}
removes 26.97\% of the \textsc{Never} baseline's wall time and \sysname
removes 34.49\%, eliminating a further 11.55\% of the stall \textsc{Lru}
leaves behind (Figure~\ref{fig:stall-ladder}). The margin widens as the budget
admits more of the working set: at 560~GB with two slots \sysname removes
50.69\% of \textsc{Lru}'s residual stall.
Figure~\ref{fig:budget-sweep} shows the effect across the full budget range,
together with the fraction of decisions in which memory pressure actually
forced an eviction. That fraction is the more informative axis: it falls from
35.2\% at 200~GB to 0.0\% by 720~GB, and where it reaches zero the correct
policy is to retain everything and no arbitration is required. Budget also
buys performance in steps rather than smoothly, since additional memory is
only useful when it crosses a threshold admitting one more resource---which is
why 320~GB and 400~GB perform almost identically.

\begin{figure}[t]
\centering
\figspec{4.2cm}{\textbf{fig:stall-ladder} --- grouped bars, stall seconds per
arm (\textsc{Never}, \textsc{Lru}, \sysname) at 256~GB/1 slot and
560~GB/2 slots. Annotate each \sysname bar with its reduction over
\textsc{Lru}.}
\caption{Stall attributable to resource residency, by policy and
configuration. \sysname removes 11.55\% of the stall \textsc{Lru} leaves at
the production allocation and 50.69\% at a larger budget.}
\label{fig:stall-ladder}
\end{figure}

\subsection{End-to-End Runtime}

\sysname reduces end-to-end wall time by 34.49\% against \textsc{Never} and
10.30\% against \textsc{Lru} at 256~GB with one device slot, rising to 50.81\%
and 17.09\% respectively with two (Figure~\ref{fig:budget-sweep}).

These wall-clock figures are a function of how much real computation the
workload performs between tool invocations, and we report that dependence
rather than a single number (Figure~\ref{fig:compute-sweep}). Holding the
mechanism fixed and varying only the inter-need gap, the wall-clock reduction
over \textsc{Lru} falls from 10.56\% at 5.7\% compute to 6.99\% at 54.9\% and
0.89\% at 95.1\%, purely because the denominator grows. The stall reduction
moves the opposite way over the same sweep---11.21\%, 15.49\%,
18.24\%---because a larger compute window hides more of each staging
operation. Our measured workload sits near the low-compute end, spending
280.3~s of a 5288.5~s trial on agent work, so we attach the compute share to
every wall-clock claim.

\begin{figure*}[t]
\centering
\figspec[0.95\textwidth]{5.0cm}{\textbf{fig:budget-sweep} --- two panels, one
per device-slot count. $x$ = budget (200--838~GB). Lines: \textsc{Lru} and
\sysname, both as \% reduction vs \textsc{Never}. Twin right-hand axis:
binding \% (share of decisions that forced an eviction), shaded. Mark the
packing thresholds at 283, 445, 562 and 838~GB with vertical rules --- the
step structure is the point.}
\caption{Wall-time reduction versus budget and device capacity. The shaded
series is the fraction of decisions in which memory pressure forced an
eviction; where it reaches zero the arbitration problem has disappeared and
retaining everything is correct.}
\label{fig:budget-sweep}
\end{figure*}

\begin{figure}[t]
\centering
\figspec{4.2cm}{\textbf{fig:compute-sweep} --- $x$ = compute share of wall
(5.7\%--95.1\%, log). Two lines vs \textsc{Lru}: wall-clock reduction
(falling) and stall reduction (rising). Mark the measured workload's operating
point near 5\%. The crossing behaviour is the reading.}
\caption{The same mechanism reported two ways. Wall-clock reduction is diluted
by compute; stall reduction improves with it, because a longer window hides
more of a staging operation.}
\label{fig:compute-sweep}
\end{figure}

\subsection{Comparison with Byte-Oriented Tiering}
\label{sec:dmsh}

Unlike the rest of this section, the results here are measured end-to-end on
hardware rather than simulated.

We compare against a state-of-the-art DMSH system on a workflow selected for
contrast: its inter-step gaps are short, so there is structurally very little
opportunity to overlap staging with computation. The DMSH arms ran only on
L40S nodes, so every row of Table~\ref{tab:dmsh} is restricted to that
partition. Pooling accelerator generations is not admissible here---identical
work has been measured to differ by up to $2.3\times$ across node types---and
the pooled figure would overstate the contrast.

\begin{table}[t]
\centering\small
\begin{tabular}{@{}lrrrl@{}}
\toprule
Configuration & $n$ & Wall (s) & $\times$base & Dominant stall class \\
\midrule
Baseline            & 9 & $270.8 \pm 52.3$  & 1.00 & no prefetch \\
\sysname            & 9 & $300.4 \pm 61.8$  & 1.11 & \texttt{no\_window} \\
DMSH, informed      & 7 & $862.2 \pm 230.9$ & 3.18 & \texttt{window\_too\_small} \\
DMSH, random        & 7 & $796.4 \pm 90.0$  & 2.94 & \texttt{window\_too\_small} \\
\bottomrule
\end{tabular}
\caption{Measured end-to-end wall time on L40S nodes for the short-window
workload. The DMSH system's informed and random staging variants are
statistically indistinguishable from one another.}
\label{tab:dmsh}
\end{table}

The DMSH system increases end-to-end wall time by $3.18\times$, and the
increase is unambiguous ($t = +6.65$). More informative than the magnitude is
the reason. Every stall it incurs classifies as \texttt{window\_too\_small}:
the system identified data to stage and began staging it, but the interval
before the need was too short to complete the transfer. Consistent with that
diagnosis, replacing its staging decisions with random ones does not make it
worse---the two variants are statistically indistinguishable ($t = +0.70$).
The failure is therefore not mis-prediction, which better prediction would
address, but the absence of a window, which it would not.

The same workload bounds what any staging-based approach can achieve on it:
an oracle arm with perfect hindsight still classifies 100\% of its stalls as
\texttt{no\_window}. \sysname is not statistically separable from the
unmodified baseline here ($t = +1.10$, on a point estimate 11\% higher), which
is the intended behaviour---where there is no window to exploit and little
reuse to retain, the system should cost approximately nothing rather than pay
for speculation that cannot repay it. We report this workload precisely
because it is the adverse case for our own mechanism as well as for the
baseline we compare against.

\subsection{Ablations}

We decompose the result by adding one mechanism at a time
(Figure~\ref{fig:ablation}). At 256~GB with one device slot, parking models
alone accounts for 5.64\%; adding retention of activated data reaches 26.97\%;
cost-aware arbitration adds 5.90 points to 32.87\%; and slack staging adds a
further 1.62 to 34.49\%. Because marginal credit depends on the order in which
mechanisms are added, we also report an order-independent Shapley
decomposition of the retention total: at this budget model parking accounts
for 2\% and data retention for 98\%, since neither of the two largest models
fits the budget at all. At 560~GB the same decomposition gives 61\% and 39\%,
so the split is a property of the allocation rather than of the system, and we
report it with the budget attached.

A second ablation varies what a wrong staging decision is permitted to destroy
(Figure~\ref{fig:prefetch-variants}). At measured predictor accuracy,
restricting staging to slack yields $+10.30$\% over \textsc{Lru} while
permitting displacement yields $+5.91$\%; restricting scope to data artifacts
yields $+13.10$\%. The safety rule matters more than the scope restriction,
and neither is a better predictor---both simply bound the cost of being wrong.

\begin{figure}[t]
\centering
\figspec{4.6cm}{\textbf{fig:ablation} --- horizontal cumulative bar,
\textsc{Never} $\rightarrow$ +model parking $\rightarrow$ +data retention
$\rightarrow$ +cost-aware arbitration $\rightarrow$ +slack staging, annotated
with each marginal contribution. Inset or second panel: Shapley split of the
retention total at 256~GB and 560~GB, showing the reversal.}
\caption{Contribution of each mechanism at 256~GB with one device slot, and
the order-independent split of the retention total at two budgets.}
\label{fig:ablation}
\end{figure}

\begin{figure}[t]
\centering
\figspec{4.2cm}{\textbf{fig:prefetch-variants} --- grouped bars at measured
accuracy: retain-only, data+slack, data+outbid, all+slack, all+outbid, as \%
vs \textsc{Lru}. Optionally repeat at two accuracies to show the safety rule's
stability.}
\caption{Two independent restrictions on speculative staging: which classes it
may stage, and whether it may displace a retained resource.}
\label{fig:prefetch-variants}
\end{figure}

\subsection{Performance Penalties of Processes on CPU for Foreground Work}

Retention and staging both consume CPU in resident worker processes, so we
measured whether that work degrades the foreground agent
(Figure~\ref{fig:cpu-interference}). With 0, 1, 2, 4 and 8 concurrent
background construction processes, foreground median latency was 0.618, 0.618,
0.617, 0.618 and 0.618~s---slowdowns of 1.000, 0.999, 0.997, 1.000 and 1.000.

A perfectly flat curve is also the signature of a probe that silently did
nothing, so we instrumented the run rather than accepting it: eight workers
were alive, load average reached 7.25, and the workers consumed 46.41
CPU-seconds in 6~s of wall time. They genuinely ran.

The result follows from construction being single-threaded, which we confirm
directly---a full artifact construction consumes 0.9885 CPU-seconds per second
of wall time. Each staging task is therefore exactly one core, and staging
tasks do not contend with one another. The practical consequence is that the
arbitrator budgets memory only, which is what makes value per gigabyte a
complete decision variable rather than one term of a two-resource trade-off.

This holds where tested, and we state the boundary. Every level measured was
under-subscribed---at most nine active processes on twelve cores---with
working sets of roughly 78~MB. Memory-bandwidth contention at the scale of a
fully activated artifact, three orders of magnitude larger, is not exercised
by this experiment.

\begin{figure}[t]
\centering
\figspec{4.2cm}{\textbf{fig:cpu-interference} --- $x$ = concurrent background
construction processes (0,1,2,4,8). Left axis: foreground median latency,
flat. Right axis: CPU-seconds consumed by the workers, rising. Plot both ---
the flatness is only credible next to evidence the workers ran.}
\caption{Foreground latency under concurrent background construction. The
rising CPU-seconds series is reported alongside the flat latency series
because a null result and a broken probe are otherwise indistinguishable.}
\label{fig:cpu-interference}
\end{figure}

\subsection{Sensitivity Sweeps}

\paragraph{Horizon.}
Retention is insensitive to $H$ across a wide range, varying by 0.14 points
between $H=60$~s and $H=600$~s, while staging degrades steadily once $H$
exceeds the workload's typical inter-need gap: 67.75\% at $H=60$, 62.19\% at
$H=600$, and 55.07\% at $H=3000$ (Figure~\ref{fig:h-sweep}). The asymmetry is
expected---$H$ exists to keep near-term candidates comparable on absolute
savings, and enlarging it past the window in which a staging operation can
complete admits candidates that cannot pay off. Setting $H$ near the observed
inter-step gap is sufficient; no tuning beyond that order of magnitude was
required.

\paragraph{Hardware topology.}
We sweep device slots against budget (Figure~\ref{fig:budget-sweep}). A second
device slot substantially raises absolute performance but reduces the value of
host-memory arbitration, because the models worth retaining are already
resident on the device and what competes for host memory is the less-used
remainder. Arbitration therefore matters most where device capacity is scarce
relative to the number of distinct models in the workflow.

\paragraph{Predictor accuracy.}
Figure~\ref{fig:accuracy-sweep} sweeps horizon accuracy from 0.15 to 1.00.
Retention is essentially flat, moving from $+7.50$\% to $+7.99$\% over
\textsc{Lru}, since horizon error perturbs a ranking rather than forcing a
selection. Slack staging degrades gracefully, from $+5.65$\% to $+18.72$\%,
and remains positive throughout: an operation that consumes only unclaimed
memory cannot destroy a retention worth more, so its downside is bounded by
wasted bandwidth. Permitting displacement removes that bound, and the same
sweep runs from $-6.22$\% at 0.15 accuracy to $+19.67$\% at 1.00, which is the
quantitative basis for the conjunctive conditions the outbid path carries in
Section~\ref{sec:prefetch}.

\begin{figure}[t]
\centering
\figspec{4.2cm}{\textbf{fig:h-sweep} --- $x = H$ (60--3000~s, log). Two
lines: retention (flat) and retention+staging (declining). Shade the region
around the workload's median inter-need gap.}
\caption{Sensitivity to the planning horizon $H$. Retention is insensitive;
staging degrades once $H$ exceeds the window in which a stage can complete.}
\label{fig:h-sweep}
\end{figure}

\begin{figure}[t]
\centering
\figspec{4.2cm}{\textbf{fig:accuracy-sweep} --- $x$ = horizon accuracy
(0.15--1.00). Three lines vs \textsc{Lru}: retention (flat), slack staging
(rising, always $>0$), outbid staging (crossing zero near 0.35). Mark the
measured accuracy band. The zero crossing is the reading.}
\caption{Sensitivity to horizon accuracy. Retention and slack staging remain
beneficial across the range; only displacement-permitting staging is
accuracy-critical.}
\label{fig:accuracy-sweep}
\end{figure}
